\chapter{BEDC.MetaCIC is its own type theory: a framing principle}
\label{ch:bedc-is-its-own-theory}

\section{The framing correction}
\label{sec:bedc-is-its-own-theory-framing-correction}

A common misread of BEDC.MetaCIC's relation to classical CIC treats
classical CIC as the reference object and BEDC.MetaCIC as a subtractive
presentation of that reference. This inverts the relationship.
BEDC.MetaCIC is a positively defined type theory with its own
first-principles commitments. Classical CIC is a different type theory
that shares some syntactic shape but adopts different design decisions.

The principle of this chapter is therefore simple:
$$
\text{BEDC.MetaCIC is its own theory, not a derivative of classical CIC.}
$$
The comparison with classical CIC is useful only after this identity has
already been fixed. It is an analytical comparison between neighbouring
theories, not a definition of BEDC.MetaCIC by reference to a parent
calculus.

\section{What BEDC.MetaCIC positively is}
\label{sec:bedc-is-its-own-theory-positive-definition}

BEDC.MetaCIC is determined by five first-principles commitments. These
commitments give the substrate its identity directly.

\subsection{Closed generating types}
\label{subsec:bedc-is-its-own-theory-closed-generators}

The term-level and certificate-level objects of the substrate come from
a finite collection of \texttt{inductive} declarations, including
\texttt{Term}, \texttt{BMark}, \texttt{BHist}, and related closed
carriers. Each declaration has a displayed constructor list. No
constructor is added by axiom, ambient library supply, quotient
identification, or unstated semantic interpretation.

The closed generator is the first unit of the theory. If an object is
not produced by the displayed constructors, then it is not an internal
object of that generator. External mathematical intentions may motivate
the naming, but they do not enlarge the internal carrier.

\subsection{Literal AST type equality}
\label{subsec:bedc-is-its-own-theory-literal-ast-equality}

Type equality on \texttt{Term}, and on any type-level object represented
inside the substrate, is literal syntactic equality of the underlying
AST. The substrate does not install a conversion rule that promotes
computational equivalence to typing identity. It does not treat two
types as identical merely because an external reduction theory relates
their encodings.

This point is central. Literal equality is not an impoverished
conversion relation. It is the equality relation chosen by the theory.
The resulting subject-reduction boundary of
\autoref{ch:metacic-subject-reduction-obstruction} is therefore an
internal fact about the chosen substrate, not a defect measured against a
different calculus.

\subsection{Structural recursors from closed generators}
\label{subsec:bedc-is-its-own-theory-structural-recursors}

For each \texttt{inductive} declaration, the Lean kernel supplies the
corresponding eliminator, such as \texttt{.rec}. These eliminators are
the universal-induction primitives available for that closed generator.
They authorize structural induction over exactly the displayed
constructors.

BEDC.MetaCIC therefore has internal induction capacity, but that
capacity is generator-local and structural. It is not a free principle
over arbitrary semantic domains. When a claim needs an induction
principle beyond the closed generator it names, that supply must be
displayed as external input rather than silently imported into the
substrate.

\subsection{Zero axiom, strict purity}
\label{subsec:bedc-is-its-own-theory-zero-axiom}

The substrate does not include \texttt{Classical.choice},
\texttt{Quot.sound}, \texttt{propext}, or any other Lean stdlib axiom as
a dependency of BEDC.MetaCIC theorems. A strict purity audit asks each
accepted theorem for its transitive axiom dependencies. For theorem
content belonging to the BEDC.MetaCIC substrate, \texttt{\#print axioms}
returns no axioms, and \texttt{axiom-purity --strict} enforces the same
boundary at project scale.

This is not a rhetorical preference for fewer assumptions. It is a
typing discipline for the project. If a result needs choice,
extensionality, quotient soundness, or propositional extensionality, the
need must appear as a named external route, a displayed setup field, or
an explicit boundary. It is not part of the substrate by default.

\subsection{Apophatic external boundary}
\label{subsec:bedc-is-its-own-theory-apophatic-boundary}

Anything beyond the preceding commitments is treated as external supply.
The supply is named apophatically: the project states what is not
internally derivable from the current substrate and exposes a socket for
that input. It does not convert the socket into a hidden axiom.

The apophatic discipline of
\autoref{ch:visions-apophatic-naming-as-discipline} is thus part of the
identity of BEDC.MetaCIC. A named hypothesis, setup field, or GAP socket
is not a disguised internal theorem. It is a visible boundary marker
where the substrate stops and an external route may begin.

\section{The defining conjunction}
\label{sec:bedc-is-its-own-theory-defining-conjunction}

The five commitments are not independent slogans. Together they define
the theory:
$$
\begin{aligned}
\mathrm{BEDC.MetaCIC}
  ={}& \text{closed generators}\\
  &+ \text{literal AST equality}\\
  &+ \text{structural recursors}\\
  &+ \text{zero-axiom strict purity}\\
  &+ \text{apophatic external boundary}.
\end{aligned}
$$
This is a positive definition. The theory is not identified by taking a
different calculus and removing a feature from it. It is identified by
the commitments it actually makes.

Consequently, the first question to ask of a BEDC.MetaCIC theorem is not
which classical theorem it resembles. The first question is which closed
generators, recursors, literal equalities, and named boundaries its proof
uses. Only after that internal reading is fixed does a comparison with
classical CIC become meaningful.

\section{Why classical CIC is a different theory}
\label{sec:bedc-is-its-own-theory-classical-cic-different}

Classical CIC has its own first-principles commitments. In familiar
presentations, a conversion rule permits typing to move across
computational equality:
$$
T_1 \equiv_{\beta} T_2
\quad\Longrightarrow\quad
(t : T_1) \Leftrightarrow (t : T_2).
$$
This rule is a type-level identification principle. It changes how
typing judgments behave and gives subject-reduction arguments access to
a kind of equality that BEDC.MetaCIC does not install.

Many classical CIC environments also provide, or commonly sit beside,
extensional and classical principles such as propositional
extensionality, function extensionality, quotient soundness, and
classical choice. Different systems choose different packages of this
supply. The shared point is that the resulting environment is organized
around a broader practical proof surface than BEDC.MetaCIC's closed
audit substrate.

The design objectives also differ. Classical CIC is built to support a
large mechanized mathematics environment. BEDC.MetaCIC is built to
delimit the smallest audit-clean closed-observational substrate on which
its meta-theory can run. These objectives produce neighbouring theories
with different commitments. Neither one is the technical parent of the
other.

\section{What comparison is good for}
\label{sec:bedc-is-its-own-theory-comparison-good}

Comparison with classical CIC is useful when it isolates a specific
design decision and its consequences. The axiom dependency map of
\autoref{ch:bedc-axiom-dependency-map} is one such comparison. It asks
how a classical-looking universal claim meets the BEDC.MetaCIC
substrate and records whether the result is an unconditional proof, a
restricted proof, or an obstruction with named supply.

In that role, comparison is secondary and diagnostic. It helps readers
see why conversion, extensionality, quotient soundness, or choice matter
for particular proof routes. It does not define BEDC.MetaCIC's identity.
The theory exists on its own footing, and its theorems are first read as
theorems of that theory.

\section{What comparison is not good for}
\label{sec:bedc-is-its-own-theory-comparison-not-good}

Subtractive slogans about BEDC.MetaCIC are structurally misleading. They
imply that classical CIC is the reference and that BEDC.MetaCIC is
valuable only as a depleted copy. That implication is false. Both
theories have first-principles definitions, and neither is the other's
reference point.

They also suggest that BEDC.MetaCIC is missing something required for
its own identity. That is false as well. BEDC.MetaCIC has exactly the
commitments that define it: closed generators, literal AST equality,
structural recursors, strict purity, and visible external boundaries. A
principle present in another calculus is not an absence inside
BEDC.MetaCIC unless BEDC.MetaCIC itself claims that principle.

Finally, they misread obstruction theorems. The subject-reduction
obstruction of \autoref{ch:metacic-subject-reduction-obstruction} is
not a failure relative to an external yardstick. It is a positive theorem
about what happens when dependent codomains are read through literal AST
equality on the BEDC.MetaCIC substrate. Classical CIC avoids that
obstruction by making different type-theoretic commitments.

\section{The label classical CIC}
\label{sec:bedc-is-its-own-theory-classical-label}

The label ``classical CIC'' is itself coarse-grained. Coq libraries,
Lean libraries, Agda libraries, and other mechanized settings make
different decisions about axioms, quotients, extensionality, universe
features, and library conventions. BEDC.MetaCIC does not target any one
of these environments as its origin.

The target is instead a specific positive substrate. Classical CIC can
serve as a comparison family, but not as a unitary reference object from
which BEDC.MetaCIC receives its meaning. Treating that family label as
the definition of BEDC.MetaCIC is a category error.

\section{Implications for related chapters}
\label{sec:bedc-is-its-own-theory-implications}

The audit map family of \autoref{ch:audit-map-family-methodology}, the
obstruction analysis of \autoref{ch:metacic-subject-reduction-obstruction},
the discharge routes of \autoref{ch:sr-discharge-routes}, the axiom
dependency map of \autoref{ch:bedc-axiom-dependency-map}, and the
closed-normal consistency result of
\autoref{ch:closed-normal-consistency-main-result} are all to be read as
internal-to-BEDC statements first.

They describe what BEDC.MetaCIC does on its own substrate. References to
classical CIC in those chapters are analytical comparisons. They are not
definitional references, and they do not turn BEDC.MetaCIC into a
subtheory relation with classical CIC.

\section{Apophatic framing of this principle}
\label{sec:bedc-is-its-own-theory-apophatic-framing}

The apophatic framing is internally consistent with this principle.
BEDC.MetaCIC names what is not derivable from its own substrate. It does
not name what another calculus possesses and then define itself by lack.

This distinction matters for the reading of every named boundary. A GAP
socket, setup field, or discharge obligation is a statement about the
current BEDC.MetaCIC substrate. It says where the internal proof route
stops. It is not a confession that the theory ought to have been a
different theory.

\section{Boundary}
\label{sec:bedc-is-its-own-theory-boundary}

This chapter is a framing principle, not a new proof and not a new Lean
target. It declares how to read BEDC.MetaCIC: as its own type theory,
defined from first principles. Other chapters may cite this principle
when their comparison language needs to keep BEDC.MetaCIC's identity
separate from the identity of classical CIC.

\concretizedIn{ch:concrete-instances-bedc-theory-identity-namecert}{BEDC theory identity is realized as a finite NameCert packet for closed generators, literal AST equality, structural recursors, strict purity, and apophatic boundary rows.}
\concretizedIn{ch:concrete-instances-metacic-theory-identity-namecert}{The MetaCIC theory-identity packet records the five positive commitments as finite NameCert rows.}
